#### `sections/2_literature_review.tex` (Обзор литературы и постановка задачи)
```latex
\section{Обзор литературы}

Основой для настоящего исследования служат работы В.А. Громова и Ф.С. Баранова  \cite{GromovBaranov2023}
, в которых предложен метод прогнозирования хаотических временных рядов,
 основанный на поиске характерных последовательностей (мотивов).
  Суть подхода заключается в реконструкции фазового пространства
   системы по наблюдаемому временному ряду и последующей кластеризации
   векторов состояний для выявления повторяющихся паттернов динамики.

Ключевой особенностью метода является использование множества мотивов,
что порождает набор возможных прогнозных значений для каждой точки.
 Для получения единого прогноза (унификации) в предыдущих работах
  рассматривались различные подходы, включая вычисление среднего
   значения, применение взвешенного усреднения (в зависимости от длины прогноза или расстояния
    до центра кластера) и выбор центра наиболее крупного кластера из возможных прогнозов.

Важнейшим аспектом, отличающим данный подход, является концепция
непрогнозируемых точек. Она позволяет алгоритму оценивать собственную уверенность
 и воздерживаться от прогноза в ситуациях высокой неопределенности,
 что повышает надежность модели. Для идентификации таких точек применялись
 как статистические критерии (энтропия множества прогнозных значений, его
 мультимодальность, ширина интерквантильного размаха), так и более сложные
 классификаторы на основе методов машинного обучения, включая логистическую регрессию,
  метод опорных векторов (SVM), деревья решений, k-NN, многослойные перцептроны (MLP) и их ансамбли.\cite{Wang2019}

Исследования показали высокую эффективность метода: алгоритм демонстрировал
практически постоянную ошибку прогноза (MAE, RMSE) даже при увеличении
горизонта до 100 шагов вперед, что достигалось за счет экспоненциального
 роста числа точек, классифицируемых как непрогнозируемые. Особое значение
  для настоящей работы имеет исследование  \cite{Gromov2021}, в котором изучалось влияние
  размера обучающей выборки на точность. Было установлено, что при увеличении о
  бъема данных из исходного ряда ошибка прогноза убывала логарифмически. Однако
  вопрос о том, можно ли компенсировать недостаток данных за счет добавления
  выборок из других, но схожих по динамике систем, оставался открытым.
 Данная работа направлена на исследование именно этой возможности.

\subsection{Формальная Постановка задачи}

Данная работа посвящена задаче многошагового прогнозирования хаотических
 временных рядов в условиях ограниченности данных. Мы рассматриваем ситуацию,
  когда доступно множество временных рядов $\mathcal{Y} = \{Y^{(0)}, Y^{(1)},
   \dots, Y^{(M)}\}$, где $Y^{(0)}$ является \textit{целевым} рядом, для которого
    необходимо построить прогноз, а $Y^{(1)}, \dots, Y^{(M)}$ — \textit{вспомогательными} рядами.
     Предполагается, что вспомогательные ряды порождены динамическими системами,
      схожими с целевой, но не обязательно идентичными ей.

Все временные ряды $Y^{(m)} = \{y_0^{(m)}, y_1^{(m)}, \dots\} \in \mathcal{Y}$ удовлетворяют следующим допущениям:
\begin{enumerate}
    \item В каждом ряду завершились все переходные процессы, и наблюдаемая
    динамика происходит в окрестности их аттрактора.
    \item Каждый ряд удовлетворяет условиям теоремы Такенса\cite{Takens1981}, что делает
    возможным корректную реконструкцию аттрактора по одномерным наблюдениям.
\end{enumerate}

Для постановки задачи прогнозирования и его последующей оценки данные разделяются
 следующим образом. Целевой ряд $Y^{(0)}$ разбивается на две непересекающиеся части:
  обучающую (наблюдаемую) выборку $Y_1^{(0)} = \{y_0^{(0)}, \dots, y_t^{(0)}\}$ и
   тестовую выборку $Y_2^{(0)} = \{y_{t+1}^{(0)}, y_{t+2}^{(0)}, \dots\}$, на которой
    производится оценка качества прогноза. Вспомогательные ряды $Y^{(1)}, \dots, Y^{(M)}$
     используются для обучения целиком. Таким образом, полное обучающее множество
     $\mathcal{D}_{\text{train}}$ определяется как объединение наблюдаемой части
      целевого ряда и всех вспомогательных рядов:
\[
\mathcal{D}_{\text{train}} = \{Y_1^{(0)}, Y^{(1)}, \dots, Y^{(M)}\}.
\]

Ключевой особенностью предлагаемого подхода является то, что для каждой
 прогнозируемой точки $y_{t+h}^{(0)}$ из тестовой выборки $Y_2^{(0)}$ алгоритм
  генерирует не одно-единственное значение, а \textit{множество возможных предсказаний}
   $\hat{S}_{t+h}$. Это множество формируется функцией $f_h$, которая использует
   как историю целевого ряда, так и всю информацию из полного обучающего множества $\mathcal{D}_{\text{train}}$:
\[
\hat{S}_{t+h} = f_h\left(\{y_t^{(0)}, y_{t-1}^{(0)}, \dots, y_{t-s+1}^{(0)}\}; \mathcal{D}_{\text{train}}\right),
\]
где $s$ — параметр, определяющий длину используемой истории.

Над множеством $\hat{S}_{t+h}$ определены два оператора. Первый оператор,
 $\zeta(\cdot)$, проверяет, является ли позиция $t+h$ прогнозируемой на
 основе анализа распределения значений в $\hat{S}_{t+h}$:
\begin{equation}
\zeta(\hat{S}_{t+h}) =
\begin{cases}
1, & \text{если позиция прогнозируема,} \\
0, & \text{в противном случае.}
\end{cases}
\end{equation}
Второй оператор, $g(\cdot)$, вычисляет единое прогнозное значение
(Unified Predicted Value, UPV) для прогнозируемой позиции:
\begin{equation}
\hat{y}_{t+h} = g(\hat{S}_{t+h}).
\end{equation}

С использованием этих операторов задача многошагового прогнозирования
 формулируется как двухкритериальная оптимизация. Необходимо одновременно
 минимизировать два функционала, вычисляемых на тестовой выборке $Y_2^{(0)}$:
\begin{enumerate}
    \item \textbf{Функционал числа непрогнозируемых точек ($I_1$):}
    \begin{equation}
    I_1 = \sum_{y_{t+h}^{(0)} \in Y_2^{(0)}} \left( 1 - \zeta(\hat{S}_{t+h}) \right) \to \min.
    \end{equation}

\item \textbf{Функционал ошибки на прогнозируемых точках ($I_2$):}
    \begin{equation}
    I_2 = \frac{1}{|Y_{\text{pred}}^{(0)}|} \sum_{y_{t+h}^{(0)} \in Y_{\text{pred}}^{(0)}} L(\hat{y}_{t+h}, y_{t+h}^{(0)}) \to \min,
    \end{equation},где $L(\hat{y}, y)$~--- функция потерь.
\end{enumerate}
где $Y_{\text{pred}}^{(0)} \subseteq Y_2^{(0)}$ — подмножество точек, для
 которых $\zeta(\hat{S}_{t+h}) = 1$. Настоящая работа направлена на решение
  этой двухкритериальной задачи путем эффективного использования расширенного
   обучающего множества $\mathcal{D}_{\text{train}}$.